\documentclass[aspectratio=169,10pt]{ctexbeamer}

\usetheme{Madrid}
\usecolortheme{default}
\setbeamertemplate{navigation symbols}{}

\usepackage{amsmath,amssymb,mathtools,bm}
\usepackage{graphicx}
\usepackage{tikz}
\usetikzlibrary{arrows.meta}


\newcommand{\e}{\mathrm{e}}

% Unified Yingxing Beamer footer and visual style.
\newcommand{\huayinglogo}{assets/logo1.jpg}
\newcommand{\huayingslogan}{英行培优，成绩无忧；英行相伴，刮目相看}

\setbeamersize{text margin left=8mm,text margin right=8mm}
\setbeamerfont{frametitle}{size=\large}
\setbeamercolor{structure}{fg=blue!70!black}
\setbeamercolor{frametitle}{bg=blue!75!black,fg=white}
\setbeamercolor{block title}{bg=blue!10,fg=black}
\setbeamercolor{block body}{bg=blue!2,fg=black}
\setbeamercolor{author in head/foot}{bg=blue!8,fg=black}
\setbeamercolor{title in head/foot}{bg=blue!8,fg=black}

\setbeamertemplate{footline}{%
  \leavevmode%
  \hbox{%
    \begin{beamercolorbox}[wd=.12\paperwidth,ht=3.8ex,dp=1.2ex,center]{author in head/foot}%
      \includegraphics[height=2.9ex]{\huayinglogo}%
    \end{beamercolorbox}%
    \begin{beamercolorbox}[wd=.76\paperwidth,ht=3.8ex,dp=1.2ex,center]{title in head/foot}%
      \scriptsize \huayingslogan
    \end{beamercolorbox}%
    \begin{beamercolorbox}[wd=.12\paperwidth,ht=3.8ex,dp=1.2ex,right]{author in head/foot}%
      \scriptsize \insertframenumber/\inserttotalframenumber\hspace*{1em}
    \end{beamercolorbox}%
  }%
}

\title{高阶导数习题课}
\subtitle{Leibniz 公式、隐函数与参数方程求导}
\author{}
\date{}

\begin{document}

\begin{frame}
  \titlepage
\end{frame}

\begin{frame}[t]{概念 1：Leibniz 公式}
\small
\begin{block}{乘积的 $n$ 阶导数}
若函数 $u(x),v(x)$ 足够光滑，则
\[
(uv)^{(n)}=\sum_{k=0}^n \binom{n}{k}u^{(k)}v^{(n-k)}.
\]
\end{block}

\begin{block}{使用要点}
\begin{enumerate}
  \item 若一个因子是低次多项式，则高阶求导后很快为 $0$；
  \item 若另一个因子是三角函数、指数函数，则高阶导数往往有规律；
  \item 因此它特别适合处理“多项式 $\times$ 基本初等函数”。
\end{enumerate}
\end{block}
\end{frame}

\begin{frame}[t]{例题 1}
\small
\begin{block}{题目}
计算
\[
\bigl(x^2\sin x\bigr)^{(80)}.
\]
\end{block}

\vspace{0.4em}
\textbf{思路提示：}用 Leibniz 公式，抓住 $x^2$ 的三阶及以上导数都为 $0$。
\end{frame}

\begin{frame}[t]{例题 1 解答}
\small
由 Leibniz 公式，
\[
\bigl(x^2\sin x\bigr)^{(80)}
=x^2(\sin x)^{(80)}
+\binom{80}{1}(x^2)'(\sin x)^{(79)}
+\binom{80}{2}(x^2)''(\sin x)^{(78)}.
\]
这里之所以只有三项，是因为
\[
(x^2)'=2x,\qquad (x^2)''=2,\qquad (x^2)^{(k)}=0\ (k\ge 3).
\]
再利用
\[
(\sin x)^{(80)}=\sin x,\quad
(\sin x)^{(79)}=-\cos x,\quad
(\sin x)^{(78)}=-\sin x.
\]
\end{frame}

\begin{frame}[t]{例题 1 解答（续）}
\small
代入得
\[
\bigl(x^2\sin x\bigr)^{(80)}
=x^2\sin x+\binom{80}{1}\cdot 2x\cdot(-\cos x)+\binom{80}{2}\cdot 2\cdot(-\sin x).
\]
又因为
\[
\binom{80}{1}=80,\qquad
\binom{80}{2}=\frac{80\cdot79}{2}=3160,
\]
所以
\[
\bigl(x^2\sin x\bigr)^{(80)}
=x^2\sin x-160x\cos x-6320\sin x.
\]
整理为
\[
\boxed{\bigl(x^2\sin x\bigr)^{(80)}=(x^2-6320)\sin x-160x\cos x.}
\]
\end{frame}

\begin{frame}[t]{例题 1 讲解提醒}
\small
\begin{block}{易错点}
\begin{enumerate}
  \item $(\sin x)^{(79)}=-\cos x,\ (\sin x)^{(78)}=-\sin x$ 的符号容易写错；
  \item 组合数 $\binom{80}{1},\binom{80}{2}$ 容易漏掉；
  \item 如果没先看出 $(x^2)^{(k)}=0\ (k\ge3)$，就会做无效展开。
\end{enumerate}
\end{block}
\end{frame}

\begin{frame}[t]{例题 2}
\small
\begin{block}{题目}
设
\[
f(x)=\frac{2x}{1-x^2},
\]
求 $f^{(n)}(x)$。
\end{block}

\vspace{0.4em}
\textbf{思路提示：}先做部分分式分解，再分别总结两类基本分式的 $n$ 阶导数。
\end{frame}

\begin{frame}[t]{例题 2 解答}
\small
先作部分分式分解。设
\[
\frac{2x}{1-x^2}=\frac{A}{1-x}+\frac{B}{1+x},
\]
则
\[
2x=A(1+x)+B(1-x)=(A+B)+(A-B)x.
\]
比较系数得
\[
A+B=0,\qquad A-B=2.
\]
解得
\[
A=1,\qquad B=-1.
\]
因此
\[
f(x)=\frac{1}{1-x}-\frac{1}{1+x}.
\]
\end{frame}

\begin{frame}[t]{例题 2 解答（续 1）}
\small
于是
\[
f^{(n)}(x)=\left(\frac{1}{1-x}\right)^{(n)}-\left(\frac{1}{1+x}\right)^{(n)}.
\]
先看第一项。因为
\[
\frac{d}{dx}(1-x)^{-1}=(-1)(1-x)^{-2}\cdot(-1)=(1-x)^{-2},
\]
继续求导可得规律
\[
\left(\frac{1}{1-x}\right)^{(n)}=\frac{n!}{(1-x)^{n+1}}.
\]
\end{frame}

\begin{frame}[t]{例题 2 解答（续 2）}
\small
再看第二项。因为
\[
\frac{d}{dx}(1+x)^{-1}=-(1+x)^{-2},
\]
继续求导可得
\[
\left(\frac{1}{1+x}\right)^{(n)}=\frac{(-1)^n n!}{(1+x)^{n+1}}.
\]
因此
\[
f^{(n)}(x)
=\frac{n!}{(1-x)^{n+1}}-\frac{(-1)^n n!}{(1+x)^{n+1}}.
\]
故
\[
\boxed{
f^{(n)}(x)=
n!\left[\frac{1}{(1-x)^{n+1}}-\frac{(-1)^n}{(1+x)^{n+1}}\right].
}
\]
\end{frame}

\begin{frame}[t]{例题 2 板书顺序}
\small
\begin{block}{板书顺序}
\begin{enumerate}
  \item 先完整写出部分分式分解；
  \item 再分别写出两类基本分式的 $n$ 阶导数；
  \item 最后合并，不要一开始就跳到结论。
\end{enumerate}
\end{block}
\end{frame}

\begin{frame}[t]{例题 3}
\small
\begin{block}{题目}
设 $y=\arcsin x$，求 $y^{(n)}(0)$。
\end{block}

\vspace{0.4em}
\textbf{思路提示：}这题不适合硬算高阶导数，要先构造递推关系，再利用奇偶性和初值。
\end{frame}

\begin{frame}[t]{例题 3 解答}
\small
先由反三角函数求导公式得
\[
y'=\frac{1}{\sqrt{1-x^2}}.
\]
改写为
\[
y'\sqrt{1-x^2}=1.
\]
两边对 $x$ 求导：
\[
y''\sqrt{1-x^2}+y'\cdot\frac12(1-x^2)^{-1/2}\cdot(-2x)=0.
\]
即
\[
y''\sqrt{1-x^2}-y'\frac{x}{\sqrt{1-x^2}}=0.
\]
两边同乘 $\sqrt{1-x^2}$，得到
\[
(1-x^2)y''-xy'=0.
\]
\end{frame}

\begin{frame}[t]{例题 3 解答（续 1）}
\small
现在对恒等式
\[
(1-x^2)y''-xy'=0
\]
作 $n$ 阶求导。

由 Leibniz 公式，
\[
\bigl[(1-x^2)y''\bigr]^{(n)}
=(1-x^2)y^{(n+2)}+n(-2x)y^{(n+1)}+\binom{n}{2}(-2)y^{(n)}.
\]
所以
\[
\bigl[(1-x^2)y''\bigr]^{(n)}
=(1-x^2)y^{(n+2)}-2nxy^{(n+1)}-n(n-1)y^{(n)}.
\]
\end{frame}

\begin{frame}[t]{例题 3 解答（续 2）}
\small
又因为
\[
(xy')^{(n)}=xy^{(n+1)}+ny^{(n)}.
\]
代回
\[
\bigl[(1-x^2)y''\bigr]^{(n)}-(xy')^{(n)}=0
\]
得到递推关系
\[
(1-x^2)y^{(n+2)}-(2n+1)xy^{(n+1)}-n^2y^{(n)}=0.
\]
令 $x=0$，得
\[
y^{(n+2)}(0)=n^2y^{(n)}(0).
\]
又因为 $\arcsin x$ 是奇函数，所以
\[
y^{(2k)}(0)=0\qquad (k=0,1,2,\dots).
\]
\end{frame}

\begin{frame}[t]{例题 3 解答（续 3）}
\small
再计算奇阶导数。由
\[
y'(0)=\frac{1}{\sqrt{1-0}}=1
\]
和递推式
\[
y^{(n+2)}(0)=n^2y^{(n)}(0),
\]
依次得
\[
y^{(3)}(0)=1^2y'(0)=1,
\]
\[
y^{(5)}(0)=3^2y^{(3)}(0)=3^2,
\]
\[
y^{(7)}(0)=5^2y^{(5)}(0)=5^2\cdot3^2.
\]
因此归纳得
\[
y^{(2k+1)}(0)=[(2k-1)!!]^2\qquad (k=0,1,2,\dots).
\]
\end{frame}

\begin{frame}[t]{例题 3 解答（续 4）}
\small
综上，
\[
\boxed{
y^{(2k)}(0)=0,\qquad
y^{(2k+1)}(0)=[(2k-1)!!]^2.
}
\]
\end{frame}

\begin{frame}[t]{例题 3 讲解提醒}
\small
\begin{block}{易错点}
\begin{enumerate}
  \item 只会算前几阶，却没有把递推式真正写出来；
  \item 忘记 $\arcsin x$ 是奇函数，从而偶阶导数在 $0$ 点全为 $0$；
  \item 把双阶乘 $(2k-1)!!$ 和普通阶乘混淆。
\end{enumerate}
\end{block}
\end{frame}

\begin{frame}[t]{例题 4}
\small
\begin{block}{题目}
设
\[
f(x)=
\begin{cases}
\e^{-1/x^2}, & x\ne 0,\\
0, & x=0,
\end{cases}
\]
证明 $f^{(n)}(0)=0\ (n\in\mathbb N_+)$。
\end{block}

\vspace{0.4em}
\textbf{思路提示：}核心不是逐阶硬算，而是抓住 $\e^{-1/x^2}$ 在 $x\to0$ 时比任何幂次都更快趋于 $0$。
\end{frame}

\begin{frame}[t]{例题 4 解答}
\small
先用一个基本事实：对任意 $\alpha>0$，
\[
\lim_{u\to+\infty}\frac{u^\alpha}{\e^u}=0.
\]
由导数定义，
\[
f'(0)=\lim_{x\to0}\frac{f(x)-f(0)}{x}
=\lim_{x\to0}\frac{\e^{-1/x^2}}{x}.
\]
令
\[
u=\frac1{x^2},
\]
则当 $x\to0$ 时，$u\to+\infty$，上式相当于
\[
u^{1/2}\e^{-u}\to0.
\]
因此
\[
f'(0)=0.
\]
\end{frame}

\begin{frame}[t]{例题 4 解答（续 1）}
\small
对 $x\ne0$，
\[
f'(x)=\left(\e^{-1/x^2}\right)'=\e^{-1/x^2}\cdot\frac{2}{x^3}
=\frac{2}{x^3}\e^{-1/x^2}.
\]
于是
\[
f''(0)=\lim_{x\to0}\frac{f'(x)-f'(0)}{x}
=\lim_{x\to0}\frac{2}{x^4}\e^{-1/x^2}.
\]
仍令
\[
u=\frac1{x^2},
\]
则上式化为
\[
2u^2\e^{-u}\to0.
\]
所以
\[
f''(0)=0.
\]
\end{frame}

\begin{frame}[t]{例题 4 解答（续 2）}
\small
继续对 $x\ne0$ 反复求导，可归纳得到
\[
f^{(n)}(x)=P_n\!\left(\frac1x\right)\e^{-1/x^2},
\]
其中 $P_n$ 是某个关于 $\dfrac1x$ 的多项式。

下面用数学归纳法证明 $f^{(n)}(0)=0$。

当 $n=1,2$ 时，上面已经证明。
设当 $n=k$ 时成立，即
\[
f^{(k)}(0)=0.
\]
则
\[
f^{(k+1)}(0)=\lim_{x\to0}\frac{f^{(k)}(x)-f^{(k)}(0)}{x}
=\lim_{x\to0}\frac1xP_k\!\left(\frac1x\right)\e^{-1/x^2}.
\]
\end{frame}

\begin{frame}[t]{例题 4 解答（续 3）}
\small
注意到
\[
\frac1xP_k\!\left(\frac1x\right)
\]
仍然只是 $\dfrac1x$ 的一个多项式，因此它的每一项都形如
\[
c_mx^{-m}.
\]
所以极限中的每一项都形如
\[
c_mx^{-m}\e^{-1/x^2}.
\]
令
\[
u=\frac1{x^2}
\]
后，就化成
\[
Cu^\beta\e^{-u}\to0.
\]
因此整个极限为 $0$，即
\[
f^{(k+1)}(0)=0.
\]
\end{frame}

\begin{frame}[t]{例题 4 解答（续 4）}
\small
由数学归纳法，对任意正整数 $n$ 都有
\[
\boxed{f^{(n)}(0)=0.}
\]
\end{frame}

\begin{frame}[t]{例题 4 讲解提醒}
\small
\begin{block}{教学提醒}
这题的本质不是“每次求导都刚好能算出来”，而是因为
\[
\e^{-1/x^2}
\]
在 $x\to0$ 时衰减得比任何有限次幂都快，所以所有这类极限都会回到 $0$。
\end{block}
\end{frame}

\begin{frame}[t]{例题 5}
\small
\begin{block}{题目}
将 $(-\infty,0]$ 上的常值函数 $u(x)\equiv0$ 和 $[1,+\infty)$ 上的常值函数 $v(x)\equiv1$ 延拓为 $(-\infty,+\infty)$ 上的无穷次可微函数，并使其值域为 $[0,1]$。
\end{block}

\vspace{0.4em}
\textbf{思路提示：}先造一个在左侧恒为 $0$、右侧平滑抬起的函数，再通过比值把两端平滑连接。
\end{frame}

\begin{frame}[t]{例题 5 解答}
\small
先定义
\[
g(x)=
\begin{cases}
0, & x\le0,\\
\e^{-1/x^2}, & x>0.
\end{cases}
\]
再令
\[
f(x)=\frac{g(x)}{g(x)+g(1-x)}.
\]
下面分区间验证它满足要求。
\end{frame}

\begin{frame}[t]{例题 5 解答（续）}
\small
\begin{enumerate}
  \item 当 $x\le0$ 时，$g(x)=0$，而 $1-x\ge1$，故 $g(1-x)>0$，于是
  \[
  f(x)=\frac{0}{0+g(1-x)}=0.
  \]
  \item 当 $x\ge1$ 时，$1-x\le0$，故 $g(1-x)=0$，且 $g(x)>0$，于是
  \[
  f(x)=\frac{g(x)}{g(x)+0}=1.
  \]
\end{enumerate}
\end{frame}

\begin{frame}[t]{例题 5 解答（续）}
\small
\begin{enumerate}
\setcounter{enumi}{2}
  \item 当 $0<x<1$ 时，$g(x)>0,\ g(1-x)>0$，所以
  \[
  0<\frac{g(x)}{g(x)+g(1-x)}<1.
  \]
  \item 由上一题知 $g$ 在整个实轴上是 $C^\infty$ 的；又因为分母恒正，所以 $f$ 也是 $C^\infty$ 的。
\end{enumerate}
因此所求延拓函数可以取为
\[
\boxed{f(x)=\frac{g(x)}{g(x)+g(1-x)}.}
\]
\end{frame}

\begin{frame}[t]{例题 5 讲解提醒}
\small
\begin{block}{易错点}
\begin{enumerate}
  \item 只验证端点常值，没有验证中间区间 $0<f(x)<1$；
  \item 忽略分母恒正这一点；
  \item 说“显然光滑”，却没有说明这是由上一题推出的。
\end{enumerate}
\end{block}
\end{frame}

\begin{frame}[t]{概念 2：隐函数求导法}
\small
\begin{block}{基本思想}
若函数关系由方程
\[
F(x,y)=0
\]
确定，并且其中 $y=y(x)$，则把 $y$ 看成 $x$ 的函数，对恒等式
\[
F(x,y(x))\equiv0
\]
两边对 $x$ 求导，从而求出 $y'(x)$。
\end{block}
\end{frame}

\begin{frame}[t]{概念 2：隐函数求导法}
\small
\begin{block}{反函数求导公式}
若 $y=y(x)$ 与 $x=x(y)$ 互为反函数，则
\[
\boxed{y'(x)=\frac{1}{x'(y)}}.
\]
\end{block}
\end{frame}

\begin{frame}[t]{例题 6}
\small
\begin{block}{题目}
设函数 $x=x(y)$ 满足 Kepler 方程
\[
y=x-q\sin x\qquad (0<q<1),
\]
求 $x'(y)$。
\end{block}

\vspace{0.4em}
\textbf{思路提示：}把 $x$ 看成 $y$ 的函数，直接对恒等式关于 $y$ 求导。
\end{frame}

\begin{frame}[t]{例题 6 解答}
\small
把 $x=x(y)$ 代入原方程，得到关于 $y$ 的恒等式
\[
y-x(y)+q\sin x(y)=0.
\]
对 $y$ 求导：
\[
1-x'(y)+q\cos x(y)\cdot x'(y)=0.
\]
把含 $x'(y)$ 的项并到一起：
\[
\bigl[-1+q\cos x(y)\bigr]x'(y)=-1.
\]
因此
\[
\boxed{x'(y)=\frac{1}{1-q\cos x(y)}.}
\]
\end{frame}

\begin{frame}[t]{例题 6 讲解提醒}
\small
\begin{block}{易错点}
\begin{enumerate}
  \item 对 $\sin x(y)$ 求导时漏写 $x'(y)$；
  \item 最终答案里把 $x(y)$ 误写成 $x$；
  \item 没说明这里是对变量 $y$ 求导。
\end{enumerate}
\end{block}
\end{frame}

\begin{frame}[t]{例题 7}
\small
\begin{block}{题目}
设函数 $y=y(x)$ 满足单位圆方程
\[
x^2+y^2=1,
\]
分别用显式方法和隐式方法求 $y'(x)$，并作比较。
\end{block}

\vspace{0.4em}
\textbf{思路提示：}显式法针对某一支，隐式法给出统一公式。
\end{frame}

\begin{frame}[t]{例题 7 解答}
\small
\textbf{显式法：}
由单位圆方程可写成两支
\[
y_1=\sqrt{1-x^2},\qquad y_2=-\sqrt{1-x^2}.
\]
分别求导，得
\[
y_1'=\frac12(1-x^2)^{-1/2}\cdot(-2x)=-\frac{x}{\sqrt{1-x^2}},
\]
\[
y_2'=-\frac12(1-x^2)^{-1/2}\cdot(-2x)=\frac{x}{\sqrt{1-x^2}}.
\]
\end{frame}

\begin{frame}[t]{例题 7 解答（续）}
\small
\textbf{隐式法：}
把 $y$ 看成 $x$ 的函数，由
\[
x^2+y^2(x)=1
\]
两边对 $x$ 求导：
\[
2x+2y(x)y'(x)=0.
\]
解得
\[
\boxed{y'(x)=-\frac{x}{y(x)}.}
\]
\end{frame}

\begin{frame}[t]{例题 7 解答（续）}
\small
当
\[
y(x)=\sqrt{1-x^2}
\]
时，代入得
\[
y'(x)=-\frac{x}{\sqrt{1-x^2}}.
\]
当
\[
y(x)=-\sqrt{1-x^2}
\]
时，代入得
\[
y'(x)=\frac{x}{\sqrt{1-x^2}}.
\]
所以隐式法给出统一公式，显式法则是分别对每一支求导。
\end{frame}

\begin{frame}[t]{例题 7 讲解提醒}
\small
\begin{block}{教学提醒}
讲这题时要提醒学生：由方程确定的“函数”不一定只有一支，隐函数公式中的 $y(x)$ 必须连同具体分支一起理解。
\end{block}
\end{frame}

\begin{frame}[t]{例题 8}
\small
\begin{block}{题目}
设函数 $y=y(x)$ 可导，且满足
\[
x^2+xy+y^2=1,
\]
求 $y'$, $y''$, $y'''$。
\end{block}

\vspace{0.4em}
\textbf{思路提示：}一直在隐式关系上求导，不要先把 $y$ 解出来。
\end{frame}

\begin{frame}[t]{例题 8 解答}
\small
先对
\[
x^2+xy+y^2=1
\]
求导，得
\[
2x+y+xy'+2yy'=0.
\]
把 $y'$ 的项合并：
\[
(x+2y)y'=-(2x+y),
\]
因此
\[
\boxed{y'=-\frac{2x+y}{x+2y}.}
\]
\end{frame}

\begin{frame}[t]{例题 8 解答（续 1）}
\small
再对恒等式
\[
2x+y+xy'+2yy'=0
\]
求导：
\[
2+y'+(xy')'+2(y')^2+2yy''=0.
\]
又因为
\[
(xy')'=y'+xy'',
\]
所以
\[
2+2y'+xy''+2(y')^2+2yy''=0.
\]
整理得
\[
(x+2y)y''+2+2y'+2(y')^2=0.
\]
\end{frame}

\begin{frame}[t]{例题 8 解答（续 2）}
\small
把
\[
y'=-\frac{2x+y}{x+2y}
\]
代入，得
\[
(x+2y)y''
=-2+2\frac{2x+y}{x+2y}-2\left(\frac{2x+y}{x+2y}\right)^2.
\]
通分后
\[
(x+2y)y''
=\frac{-2(x+2y)^2+2(2x+y)(x+2y)-2(2x+y)^2}{(x+2y)^2}.
\]
\end{frame}

\begin{frame}[t]{例题 8 解答（续 3）}
\small
展开并整理分子：
\[
-2(x^2+4xy+4y^2)+2(2x^2+5xy+2y^2)-2(4x^2+4xy+y^2)
\]
\[
=-6x^2-6xy-6y^2
\]
\[
=-6(x^2+xy+y^2).
\]
而题设给出
\[
x^2+xy+y^2=1,
\]
故
\[
(x+2y)y''=-\frac{6}{(x+2y)^2}.
\]
\end{frame}

\begin{frame}[t]{例题 8 解答（续 4）}
\small
于是
\[
\boxed{y''=-\frac{6}{(x+2y)^3}.}
\]

再对上式求导：
\[
y'''=-6\cdot(-3)(x+2y)^{-4}(1+2y').
\]
又由
\[
1+2y'=1-\frac{2(2x+y)}{x+2y}
=\frac{x+2y-4x-2y}{x+2y}
=-\frac{3x}{x+2y},
\]
得
\[
\boxed{y'''=-\frac{54x}{(x+2y)^5}.}
\]
\end{frame}

\begin{frame}[t]{例题 8 讲解提醒}
\small
\begin{block}{易错点}
\begin{enumerate}
  \item 二次求导时漏掉 $(xy')'$ 或 $2yy''$；
  \item 代入 $y'$ 后通分化简不彻底；
  \item 最后一步忘了再代回 $1+2y'$。
\end{enumerate}
\end{block}
\end{frame}

\begin{frame}[t]{概念 3：参数方程求导法}
\small
\begin{block}{参数方程}
若
\[
x=\varphi(t),\qquad y=\psi(t),
\]
并且 $\varphi'(t)\ne0$，则在局部可把 $y$ 看成 $x$ 的函数，并有
\[
\boxed{\frac{dy}{dx}=\frac{\psi'(t)}{\varphi'(t)}}.
\]
\end{block}
\end{frame}

\begin{frame}[t]{概念 3：参数方程求导法}
\small
\begin{block}{继续求高阶导数}
先把 $\dfrac{dy}{dx}$ 看成关于参数 $t$ 的函数，再对 $x$ 求导：
\[
\frac{d^2y}{dx^2}
=\frac{d}{dt}\!\left(\frac{dy}{dx}\right)\Big/\frac{dx}{dt},
\]
\[
\frac{d^3y}{dx^3}
=\frac{d}{dt}\!\left(\frac{d^2y}{dx^2}\right)\Big/\frac{dx}{dt}.
\]
\end{block}
\end{frame}

\begin{frame}[t]{例题 9}
\small
\begin{block}{题目}
将例题 8 中的椭圆用参数方程表示为
\[
x=\frac{2}{\sqrt3}\cos t,\qquad
y=\sin t-\frac{1}{\sqrt3}\cos t,
\]
求 $y$ 对 $x$ 的前三阶导数：$y_x'$, $y_{xx}''$, $y_{xxx}'''$。
\end{block}

\vspace{0.4em}
\textbf{思路提示：}先算 $x_t',y_t'$，再求 $\dfrac{dy}{dx}$，之后统一对 $t$ 求导再除以 $x_t'$。
\end{frame}

\begin{frame}[t]{例题 9 解答}
\small
先求参数导数：
\[
x_t'=-\frac{2}{\sqrt3}\sin t,\qquad
y_t'=\cos t+\frac{1}{\sqrt3}\sin t.
\]
因此
\[
y_x'=\frac{y_t'}{x_t'}
=\frac{\cos t+\frac1{\sqrt3}\sin t}{-\frac{2}{\sqrt3}\sin t}.
\]
把分子分母同乘 $\sqrt3$，得
\[
y_x'=\frac{\sqrt3\cos t+\sin t}{-2\sin t}
=-\frac{\sqrt3}{2}\cot t-\frac12.
\]
\end{frame}

\begin{frame}[t]{例题 9 解答（续 1）}
\small
再求二阶导数：
\[
y_{xx}''=\frac{d}{dt}(y_x')\Big/x_t'.
\]
先算
\[
\frac{d}{dt}\left(-\frac{\sqrt3}{2}\cot t-\frac12\right)
=\frac{\sqrt3}{2}\csc^2 t.
\]
所以
\[
y_{xx}''
=\frac{\frac{\sqrt3}{2}\csc^2 t}{-\frac{2}{\sqrt3}\sin t}
=-\frac34\csc^3 t.
\]
\end{frame}

\begin{frame}[t]{例题 9 解答（续 2）}
\small
最后求三阶导数：
\[
y_{xxx}'''=\frac{d}{dt}(y_{xx}'')\Big/x_t'.
\]
因为
\[
\frac{d}{dt}\left(-\frac34\csc^3 t\right)
=-\frac34\cdot 3\csc^2 t\cdot(-\csc t\cot t)
=\frac94\csc^3 t\cot t,
\]
所以
\[
y_{xxx}'''
=\frac{\frac94\csc^3 t\cot t}{-\frac{2}{\sqrt3}\sin t}
=-\frac{9\sqrt3}{8}\frac{\cos t}{\sin^5 t}.
\]
\end{frame}

\begin{frame}[t]{例题 9 解答（续 3）}
\small
因此
\[
\boxed{
y_x'=-\frac{\sqrt3}{2}\cot t-\frac12,\quad
y_{xx}''=-\frac34\csc^3 t,\quad
y_{xxx}'''=-\frac{9\sqrt3}{8}\frac{\cos t}{\sin^5 t}.
}
\]
\end{frame}

\begin{frame}[t]{例题 9 板书顺序}
\small
\begin{block}{板书顺序}
\begin{enumerate}
  \item 先写 $x_t',y_t'$；
  \item 再写 $y_x'=\dfrac{y_t'}{x_t'}$；
  \item 之后统一用“对 $t$ 求导再除以 $x_t'$”来推二阶、三阶导数。
\end{enumerate}
\end{block}
\end{frame}

\begin{frame}[t]{小结}
\small
\begin{block}{本节三条主线}
\begin{enumerate}
  \item Leibniz 公式：适合处理乘积型高阶导数；
  \item 隐函数求导：把方程当作恒等式直接求导；
  \item 参数方程求导：先对参数求导，再除以 $x_t'$。
\end{enumerate}
\end{block}

\begin{alertblock}{做题习惯}
高阶导数题不要盲目硬算，先找结构、找规律、找递推；该展开的中间步骤要写清楚，尤其是求导来源和代入化简过程。
\end{alertblock}
\end{frame}

\end{document}

